% ============================================================
% 4. IMPLEMENTACIÓN
% (preparación del sistema, gestión de usuarios, gestión de
%  demonios con systemd y balanceo de carga con HAProxy)
% ============================================================
\section{Implementación}
\label{sec:implementacion}

La capa HTTP está implementada: los catorce recursos de la
\cref{tab:api-rest} tienen su manejador, la validación del servidor
aplica las reglas descritas en la \cref{sec:contrato-rest} y cada
petición se responde con el modelo JSON de la
\cref{sec:vision-general}. El descubrimiento de sesión
expone \texttt{GET /api/session} y el módulo de paquetes ya ejecuta APT
sobre la máquina. Los demás casos de uso, mientras no se conecten,
responden \texttt{501} con un mensaje comprensible, de modo que la
interfaz puede integrarse sin esperar a la ejecución real. Las pruebas
cubren el 96.0\,\% de los manejadores y el 100\,\% de los servicios.

% ------------------------------------------------------------
% 4.1 Preparación del sistema
% ------------------------------------------------------------
\subsection{Preparación del Sistema}
\label{sec:preparacion}

\subsubsection{Instalación y Eliminación de Aplicaciones}
\label{sec:apt}

El módulo de paquetes se apoya en el recurso \texttt{/api/packages}. El
formulario indica la acción y el nombre del paquete; el servidor valida
el nombre con la regla de paquetes APT antes de tocar la máquina. La
instalación ejecuta \texttt{apt-get install -y} y la eliminación
\texttt{apt-get remove -y} sobre ese nombre, con
\texttt{DEBIAN\_FRONTEND=noninteractive} para que APT no solicite
confirmaciones aunque la sesión no sea interactiva.

Cada operación abre una sesión con el conector disponible, la cierra al
terminar y clasifica el resultado. Si el código de salida es cero,
responde \texttt{201} al instalar (con \texttt{Location}) o
\texttt{204} al eliminar, e incluye la salida relevante; si es distinto
de cero, responde \texttt{500} conservando la orden ejecutada, el
código de salida y el \texttt{stderr} para el diagnóstico. El listado
usa \texttt{dpkg-query} y filtra por el estado
\texttt{install ok installed}. La cobertura del servicio de paquetes es
del 100\,\%.

\subsubsection{Configuración de la Llave de Acceso SSH para el Usuario Root}
\label{sec:ssh-root}
% TODO: creación de llaves para root, instalación de la llave pública en
% authorized_keys, deshabilitación del acceso por login y password,
% descarga de la llave privada y mensaje de recomendación de uso
% correcto de la llave.

\subsubsection{Personalización de la Entrada del Sistema}
\label{sec:issue}
% TODO: configuración de /etc/issue con opciones para nombre del host,
% fecha, hora, dirección IP y mensaje personalizado de inicio de sesión.

\subsubsection{Asignación del Nombre de Host y Actualización del Archivo de Hosts}
\label{sec:hostname}
% TODO: establecimiento del nombre de host y actualización del
% contenido de /etc/hosts.

% ------------------------------------------------------------
% 4.2 Gestión de usuarios
% ------------------------------------------------------------
\subsection{Gestión de Usuarios}
\label{sec:usuarios}

\subsubsection{Creación de Usuarios y Privilegios de sudo}
\label{sec:crear-usuario}
% TODO: creación de usuario con casilla de verificación para privilegios
% sudo; casilla para generar el par de llaves criptográficas, instalar
% la llave pública y habilitar el descargue de la llave privada;
% deshabilitación del acceso por login y password al habilitar SSH.

\subsubsection{Eliminación de Usuarios}
\label{sec:eliminar-usuario}
% TODO: eliminación de usuarios con casilla de verificación que indica
% si también deben eliminarse todos sus archivos.

% ------------------------------------------------------------
% 4.3 Gestión de demonios con systemd
% ------------------------------------------------------------
\subsection{Gestión de Demonios con systemd}
\label{sec:systemd}
% TODO: configuración de una aplicación como servicio de systemd
% (por ejemplo, la aplicación web del catálogo de libros): transferencia
% de los archivos requeridos a la máquina virtual y gestión de
% permisos; botones start, stop, restart, status, enable y disable;
% prueba complementaria con el script que agrega fecha y hora a un
% archivo.

% ------------------------------------------------------------
% 4.4 Balanceo de carga con HAProxy
% ------------------------------------------------------------
\subsection{Gestión de un Balanceador de Carga con HAProxy}
\label{sec:haproxy}
% TODO: mecanismo para agregar y eliminar un servidor del back-end del
% balanceador de carga; verificación del balanceo con back-end creadas
% manualmente con VirtualBox (entorno de pruebas, fuera de la consola).
